\documentclass[10pt]{extarticle}

% ---------- Packages ----------
\usepackage{fancyhdr}
\usepackage{geometry}
\usepackage{hyperref}
\usepackage{xcolor}
\usepackage{amsmath, amsfonts, bm}
\usepackage{subcaption}
\usepackage{graphicx}
\usepackage{tikz}
\usepackage{nomencl}
\usepackage{array}
\usepackage{pdflscape}
\usepackage[export]{adjustbox}
\usepackage[utf8]{inputenc}
\usepackage{parskip}

\usetikzlibrary{shapes.geometric}
\usetikzlibrary{arrows.meta, arrows}

% ---------- Page setup ----------
\geometry{
    a4paper,
    total={170mm,257mm},
    left=20mm,
    top=20mm
}

% ---------- Headers and footers ----------
\pagestyle{fancy}
\pagenumbering{arabic}

\fancyhead[L]{Rocket Physics by Yaroslav Knyazkov}
\fancyfoot[C]{}
\fancyfoot[R]{\thepage}

% ---------- Hyperlinks ----------
\hypersetup{
    colorlinks=true,
    linkcolor={blue!50!black},
    citecolor={blue!50!black},
    urlcolor={blue!80!black}
}

% ---------- Equation numbering ----------
\numberwithin{equation}{subsection}

% ---------- Custom commands ----------
\newcommand{\todo}[1]{%
    \textcolor{red}{(#1)}%
}
% \ai marks text/equations written or edited by the AI assistant, so changes
% made outside a normal editing pass are visually distinguishable.
\newcommand{\ai}[1]{{\color{blue}#1}}

\begin{document}

\tableofcontents

\ai{This report is organised in two halves. \autoref{sec:rocket2d} fixes the frames, the three
forces and their application points, writes the equations of motion and analyses them by hand;
\autoref{sec:models} gives the interchangeable models behind each symbol those equations use.
\autoref{sec:assumptions}, immediately below, lists every simplification both halves currently
rely on, and \autoref{sec:backlog} lists what is deliberately deferred. Text and equations in blue
were written or edited by an AI assistant; parenthesised text in red marks an item that is still
open.}

\newpage
\ai{\section{Current assumptions}\label{sec:assumptions}
This section collects, in one place, every simplification the current derivation and the
simulator rely on. Each is expected to be relaxed in a later phase; the relevant model or roadmap
step is named where one exists, so a reader can tell ``current shortcut'' from ``permanent
choice'' at a glance.
\begin{itemize}
    \item \textbf{Flat Earth for the flight-dynamics frame.} $\hat{z}_g$ is fixed for the whole
    flight: the local vertical does not rotate with downrange distance
    (\autoref{sec:gravity} > Direction). The Earth's rotation is carried only in part. Its
    \emph{static} effect is present --- the \emph{magnitude} of $g$ includes the centrifugal
    contribution through the WGS84 model (\autoref{sec:gravity}) --- and so is the in-plane
    Coriolis term in the translational equations of motion
    (\autoref{sec:noninertial}). What is dropped is the out-of-plane Coriolis
    component, which this 2D model has no degree of freedom to receive, and the residual
    centrifugal term, derived there but left unimplemented as some three orders of magnitude
    smaller. The frozen-$\hat{z}_g$ error and the retained Coriolis term are of comparable
    magnitude; see that section's closing caveat before over-reading either.
    \item \textbf{Constant moment arms (not constant mass).} Mass, centre of gravity and moment
    of inertia are \emph{not} held fixed: $m(t)$, $l_B(t)$ and $I(t)$ come from the propellant
    model of \autoref{sec:mass} and are carried by \autoref{sec:eom}. What is still held
    fixed is one step further out --- how
    the other two application points travel relative to a moving centre of gravity. The
    CoG-to-$D$ arm $l_1 = l_D - l_B$ responds to $(\alpha, M)$ through the coefficient table
    (\autoref{sec:aero}) but not to CoG travel, and the CoG-to-$C$ thrust arm is a
    constant. Lifting this needs a stated rule for how $C$ and $D$ move as $l_B(t)$ moves, which
    no section currently gives.
    \item \textbf{Thrust axis aligned with the body axis.} No gimbal or nozzle misalignment is
    modelled; under the current geometry the thrust moment about the CoG is identically zero
    (\autoref{sec:analysis_moment}). See \autoref{sec:backlog}.
    \item \textbf{Planar (2D) flight, no crosswind.} Motion is confined to the vertical plane
    fixed by the launch azimuth $\chi$; the out-of-plane wind component is computed and discarded
    (\autoref{sec:wind}).
    \item \textbf{No vertical wind.} Wind velocity is horizontal only; updrafts/downdrafts are not
    modelled (\autoref{sec:wind}).
    \item \textbf{No launch rail constraint, no recovery/parachute.} Free flight from $t = 0$;
    see \autoref{sec:backlog}.
\end{itemize}
\autoref{sec:backlog} is the complement of this list: this one records what the present model
simplifies, that one what is deliberately left undone and what each item is waiting on.}

\newpage
\section{Rocket 2D}
\subsection{Frames}
\begin{itemize}
    \item $\mathcal{F}_0$ - ground
    \item $\mathcal{F}_g$ - gravitational field
    \item $\mathcal{F}_r$ - rocket
\end{itemize}
\ai{$\mathcal{F}_0$ and $\mathcal{F}_g$ share the same fixed orientation ($\hat{z}_g \equiv
\hat{z}_0$, down) --- the gravity frame is kept notationally separate only so a curved/rotating
Earth can be introduced later without renaming anything (\S Current assumptions). $\mathcal{F}_r$
is right-handed with $\hat{x}_r$ along the rocket's nose axis (positive forward), $\hat{z}_r$
completing the pitch plane with it ($\hat{z}_r \equiv \hat{z}_g$ when $\theta = 0$), and
$\hat{y}_r$ out of the pitch plane. Pitch $\theta$ is measured from the inertial horizontal,
positive nose-up; the pitch rate $\omega = \dot\theta$ is the rotation rate about $\hat{y}_r$.}


\subsection{Forces and moments}
Weight is the main force the rocket has to overcome to increase its altitude. It is written by
\begin{equation}
    \vec{W} = mg\hat{z}_g, 
\end{equation}
where $m$ is mass of rocket, $g$ gravitational acceleration along direction of gravitational field $\hat{z}_g$.

To overcome weight thrust force is generated, by ejecting fuel at velocity and due to pressure differences. It is simply written as
\begin{equation}
    \vec{T} = T\hat{x}_r,
\end{equation}
where $T$ is thrust magnitude along axis of rocket $\hat{x}_r$. \todo{add own thrust axis; misalignment + control}

Due to motion of rocket within fluid aerodynamic forces are present and given by
\begin{equation}
    \vec{A}_f = -A_x\hat{x}_r - A_z\hat{z}_r,
\end{equation}
where $A_x, A_z$ are respectively drag and lift magnitudes defined in rocket frame $\hat{x}_r, \hat{z}_r$. By defining a reference point $D$ on the rocket pitching moment can be obtained and written as
\begin{equation}
    \vec{A}_m = -A_m\hat{y}_r,
\end{equation}
where $A_m$ is pitching moment magnitude, defined around $\hat{y}_r$ axis.

\subsection{Geometry}
Three forces, each act at it's own point. Every dimensions will be defined in rocket frame and from the point $A$ at the tip of the rocket.

Weight is defined at center of gravity, here defined as $B$, giving
\begin{equation}
    \vec{AB} = -l_B\hat{x}_r - u_B\hat{z}_r.
\end{equation}

Thrust is defined at $C$,
\begin{equation}
    \vec{AC} = -l_C\hat{x}_r - u_C\hat{z}_r.
\end{equation}

Aerodynamic forces at $D$
\begin{equation}
    \vec{AD} = -l_D\hat{x}_r.
\end{equation}

The added $u$ dimensions refer to the axis misalignment as in real rockets center of gravity and thrust can't be strictly on the axis. As $D$ is reference point no misalignment definition is required, however the external surfaces unsymmetry will be visible in polar plots.

\subsection{Equations of motion}
\begin{gather}
    m\vec{\ddot{x}} = \sum \vec{F} = \vec{W} + \vec{T} + \vec{A}_f \\
    I\ddot{\theta} = \sum \vec{M}(B) \cdot \hat{y}_r = \left[\vec{BC}\times\vec{T} + \vec{BD}\times\vec{A}_f + \vec{A_m}\right]\cdot \hat{y}_r
\end{gather}

\subsection{Analytical analysis}
\ai{Here ${}^r(\cdot)$ denotes a vector's components resolved along the body axes $\hat{x}_r,
\hat{z}_r$ at the current instant --- i.e.\ $m\,{}^r\ddot{x} = (m\ddot{\vec{x}})\cdot\hat{x}_r$
--- \emph{not} a derivative taken within the rotating frame $\mathcal{F}_r$ itself (which would
additionally need that frame's own rotation-coupling terms). The equations below are therefore
just $\sum\vec{F}\cdot\hat{x}_r$, $\sum\vec{F}\cdot\hat{z}_r$ of \S Equations of motion,
expanded.}
\subsubsection{Forces}
\begin{align}
    m{}^r\ddot{x} &= -mg\cos{\theta} + T - A_x\\
    m{}^r\ddot{z} &= mg\sin{\theta} - A_z
\end{align}

\subsubsection{Moment}
\begin{equation}
    I\ddot{\theta} = (u_B - u_D)T - u_BA_x + (l_B - l_C)A_z - A_m
\end{equation}
\todo{write that $l_C$ is irrelevant; unless thrust axis defined}

No thrust \todo{write}
\begin{equation}
    I\ddot{\theta} = -u_BA_x + (l_B - l_C)A_z - A_m.
\end{equation}

If thrust and weight axis are aligned then the equation turns into
\begin{equation} \label{eq:moment_aligned}
    I\ddot{\theta} = (l_B - l_C)A_z - A_m.
\end{equation}

\todo{actual section with thrust and no-thrust stability}
\paragraph{Static stability} is a steady measure of stability of rocket due to disturbance, which for this case is angle of attack $\alpha$. We will use no thrust, axis aligned \autoref{eq:moment_aligned}. Before the analysis we need to expand and linearize $A_z, A_m$, which are given by
\begin{equation}
    A_z = qc_zA,\quad A_m = qc_mAl,
\end{equation}
where $q$ is dynamic pressure, $c_z, c_m$ are respectively lift and pitching moment coefficeints and $A, l$ reference area and length, More on that in \todo{section reference}.

Coefficents $c_z, c_m$ can be linearized by using
\begin{align}
    c_z &= c_{z,0} + \alpha c_{z,\alpha}, \\
    c_m &= c_{m,0} + \alpha c_{m,\alpha}.
\end{align}

By taking partial derivative of angular acceleration by $\alpha$, we expect $\frac{\partial \ddot\theta}{\partial \alpha} < 0$ behavior, meaning the rocket will stabilize the trajectory. The expansion and derivative gives
\begin{equation}
    \frac{\partial \ddot\theta}{\partial \alpha} = \frac{qA}{I}\biggl[(l_B - l_C)c_{z,\alpha} - c_{m,\alpha}l\biggr] < 0,
\end{equation}
which is met stability margin $\textrm{SM}$ is 
\begin{equation}
    \textrm{SM} = l_B - l_C < l\frac{c_{m,\alpha}}{c_{z,\alpha}}.
\end{equation}
\todo{maybe divide both sides by area; sqrtA or diameter for dimensionless number}

\newpage
\section{Models}

\subsection{Relative wind and angle of attack}
Let $\vec{v}$ be the rocket's velocity relative to $\mathcal{F}_0$, and $\vec{v}_w$ the wind
velocity relative to $\mathcal{F}_0$ (a model output, \S Wind). The rocket's velocity
relative to the surrounding air is
\begin{equation}
    \vec{v}_{rel} = \vec{v} - \vec{v}_w = u\hat{x}_r + w\hat{z}_r
\end{equation}
with $u$ the axial (along-nose) component and $w$ the normal component, both expressed in
$\mathcal{F}_r$. The relative airspeed and angle of attack are
\ai{\begin{equation}
    V = \sqrt{u^2+w^2}, \qquad \alpha = \operatorname{atan2}(w, u)
\end{equation}}
so that $u = V\cos\alpha$, $w = V\sin\alpha$. At $\alpha = 0$ (nose-first into the relative wind)
$\vec{v}_{rel} = V\hat{x}_r$.

Define the unit vector along the relative wind and its in-plane normal (a $+90^\circ$ rotation of
$\hat{u}$ within $\mathcal{F}_r$):
\begin{equation}
    \hat{u} = \frac{\vec{v}_{rel}}{V} = \cos\alpha\,\hat{x}_r + \sin\alpha\,\hat{z}_r, \qquad
    \hat{n} = -\sin\alpha\,\hat{x}_r + \cos\alpha\,\hat{z}_r
\end{equation}
At $\alpha = 0$: $\hat{u} = \hat{x}_r$, $\hat{n} = \hat{z}_r$.
\ai{$\hat{u}, \hat{n}$ feed the dynamic pressure $q$ and, through $\alpha$, the coefficient
lookup in \S Aerodynamics; the aerodynamic force itself is resolved along the fixed body axes
$\hat{x}_r, \hat{z}_r$, \emph{not} along $\hat{u}, \hat{n}$ --- see \S Aerodynamics.}

\subsection{Gravity}
Interchangeable models for the gravitational acceleration magnitude $g$ used in
$\vec{W} = mg\hat{z}_g$ (\S Forces). All variants act along the fixed $\hat{z}_g$ direction under
the flat, non-rotating Earth assumption (\S Current assumptions); $h = -z$ is altitude above the
launch point, $R_e$ the Earth's radius.

\paragraph{Magnitude.}
\begin{itemize}
    \item \textbf{Constant} (Phase 1 baseline): $g = g_0$.
    \item \textbf{Radial} (inverse-square): $g(h) = g_0\left(\dfrac{R_e}{R_e + h}\right)^2$ ---
    exact for a spherically symmetric, non-rotating Earth; still directed along $\hat{z}_g$ under
    the flat-range assumption.
    \item \textbf{Real data (WGS84)}: Somigliana normal gravity at the ellipsoid surface, plus the
    second-order free-air correction to height above it:
    \begin{gather}
        \gamma_0(\phi) = \gamma_e\dfrac{1+k\sin^2\phi}{\sqrt{1-e^2\sin^2\phi}} \\
        g(\phi, h) = \gamma_0\left[1 - \dfrac{2}{a_e}\left(1+f+m-2f\sin^2\phi\right)h + \dfrac{3}{a_e^2}h^2\right]
    \end{gather}
    with $\phi$ the launch site's geodetic latitude (fixed for the flight, \S Current assumptions),
    $h$ the altitude above the ellipsoid, and published WGS84 constants $a_e = 6378137\,\mathrm{m}$,
    $f = 1/298.257223563$, $m = 0.00344978650684$, $\gamma_e = 9.7803253359\,\mathrm{m/s^2}$,
    $k = 0.00193185265241$, $e^2 = 0.00669437999013$. A gridded dataset (e.g.\ EGM2008) would
    resolve anomalies two orders of magnitude below this model's own error and is not used here.
    \ai{This variant already includes the Earth's rotation (centrifugal) contribution to the
    \emph{magnitude} of $g$ through $\gamma_0(\phi)$ --- the equatorial-vs-polar difference in $g$
    is a rotation effect, folded into the Somigliana formula. That is a separate, static effect
    from the dynamic (Coriolis) correction to the equations of motion, \S Non-inertial frame
    corrections below.}
\end{itemize}

\paragraph{Direction.}
\ai{True gravitational acceleration points radially inward, toward the Earth's center. Over a
flight's downrange extent $x$, the true local-vertical direction rotates by approximately
$x/R_e$ relative to the launch-site vertical --- for a suborbital flight of a few kilometers to a
few tens of kilometers this is a fraction of a milliradian, negligible against every other error
source in this model. $\hat{z}_g$ is therefore frozen at the launch site's local vertical for the
whole flight (\S Current assumptions). A downrange-dependent gravity \emph{direction} would be the
one genuinely new physics among the gravity variants above, but it needs this assumption lifted
first; see the backlog.}

\subsection{Non-inertial frame corrections (Earth's rotation)}
\ai{The WGS84 gravity model above already accounts for the Earth's rotation in the magnitude of
$g$ (the static, equilibrium-shape effect, folded into $\gamma_0(\phi)$). A rotating Earth also
introduces two velocity/position-dependent fictitious forces into the translational equations of
motion once $\mathcal{F}_0$ is reinterpreted as Earth-fixed rather than inertial:
\begin{gather}
    \vec{a}_{Coriolis} = -2\vec{\Omega}\times\vec{v} \\
    \vec{a}_{centrifugal} = -\vec{\Omega}\times(\vec{\Omega}\times\vec{r})
\end{gather}
with $\vec{\Omega}$ the Earth's angular velocity ($|\vec{\Omega}| = 7.2921\times10^{-5}\,
\mathrm{rad/s}$) and $\vec{r}, \vec{v}$ the rocket's position/velocity relative to the Earth-fixed
frame. The centrifugal term is a small, latitude-dependent correction, already absorbed into the
WGS84 $g(\phi,h)$ above when the launch site is the reference point; the Coriolis term is
genuinely new here and depends on the flight's own velocity and heading, not just the launch
site.
\paragraph{Tangent-plane basis.} Work in the local East--North--Up triad
$(\hat{E},\hat{N},\hat{U})$ at the launch site $(\phi_0,\lambda_0)$, with $\phi_0$ the site's
geodetic latitude (the same quantity written $\phi$ in \S Gravity). The flight-plane frame
$\mathcal{F}_0$ is the NED-like triad already fixed by the launch heading $\chi$ (degrees
clockwise from North, the direction of $+\hat{x}$, \S Atmosphere > Field) and by $\hat{z}_g$
pointing \emph{down} (\S Frames, $h=-z$):
\begin{gather}
    \hat{x}_g = \sin\chi\,\hat{E} + \cos\chi\,\hat{N} \\
    \hat{y}_g = \cos\chi\,\hat{E} - \sin\chi\,\hat{N} \\
    \hat{z}_g = -\hat{U}
\end{gather}
Nothing here is new: $\hat{x}_g$ is the downrange direction used by the ground track
$(x\sin\chi, x\cos\chi)$ of \S Atmosphere > Field, and $\hat{y}_g$ is exactly the cross-plane
direction whose wind component $w_y = u\cos\chi - v\sin\chi$ is computed and discarded in
\S Wind > Flight-plane projection. The triad is right-handed,
$\hat{x}_g\times\hat{y}_g = \hat{z}_g$, which can be checked directly from the three lines above
($(\sin\chi,\cos\chi,0)\times(\cos\chi,-\sin\chi,0) = (0,0,-1)$ in ENU components). $\hat{y}_g$
points to the \emph{right} of the downrange direction.

\paragraph{Components of $\vec{\Omega}$.} The Earth's angular velocity points along the rotation
axis, i.e.\ toward the celestial north pole. At latitude $\phi_0$ that direction lies in the
local meridian plane (no East component), tilted from the local vertical by the co-latitude:
\begin{equation}
    \vec{\Omega} = \Omega\left(\cos\phi_0\,\hat{N} + \sin\phi_0\,\hat{U}\right), \qquad
    \Omega = 7.2921\times10^{-5}\,\mathrm{rad/s}
\end{equation}
Projecting onto the flight-plane triad by taking dot products with the three basis vectors above:
\begin{align}
    \Omega_x &= \vec{\Omega}\cdot\hat{x}_g
        = \Omega\cos\phi_0\,(\hat{N}\cdot\hat{N})\cos\chi = \Omega\cos\phi_0\cos\chi \\
    \Omega_y &= \vec{\Omega}\cdot\hat{y}_g
        = \Omega\cos\phi_0\,(\hat{N}\cdot\hat{N})(-\sin\chi) = -\Omega\cos\phi_0\sin\chi \\
    \Omega_z &= \vec{\Omega}\cdot\hat{z}_g
        = \Omega\sin\phi_0\,(\hat{U}\cdot(-\hat{U})) = -\Omega\sin\phi_0
\end{align}
so that
\begin{equation} \label{eq:omega_components}
    \vec{\Omega} = \Omega\left(\cos\phi_0\cos\chi\,\hat{x}_g
        - \cos\phi_0\sin\chi\,\hat{y}_g - \sin\phi_0\,\hat{z}_g\right)
\end{equation}
with $\Omega_x^2+\Omega_y^2+\Omega_z^2 = \Omega^2$ as required. Three sanity checks: at the North
pole ($\phi_0 = 90^\circ$) it reduces to $-\Omega\hat{z}_g$, i.e.\ straight up, for every $\chi$;
on the equator launching due North ($\phi_0=0,\chi=0$) to $+\Omega\hat{x}_g$, i.e.\ horizontal and
downrange; on the equator launching due East ($\phi_0=0,\chi=90^\circ$) to $-\Omega\hat{y}_g$,
which is horizontal and to the left of downrange --- North again, as it must be.

The sign of $\Omega_z$ is the one easy mistake here and is purely a bookkeeping consequence of
this project's $z$-down convention: the rotation vector points \emph{up} out of the northern
hemisphere, hence negative along $\hat{z}_g$.

\paragraph{Which component of $\vec{\Omega}$ actually matters in 2D.} With motion confined to
the flight plane, $\vec{v} = v_x\hat{x}_g + v_z\hat{z}_g$ (no $\hat{y}_g$ component). Expanding
the cross product in the right-handed triad:
\begin{equation}
    \vec{\Omega}\times\vec{v} =
        \Omega_yv_z\,\hat{x}_g
        + \left(\Omega_zv_x - \Omega_xv_z\right)\hat{y}_g
        - \Omega_yv_x\,\hat{z}_g
\end{equation}
so
\begin{equation} \label{eq:coriolis_split}
    \vec{a}_{Coriolis} = -2\vec{\Omega}\times\vec{v} =
        \underbrace{-2\Omega_y\left(v_z\hat{x}_g - v_x\hat{z}_g\right)}_{\text{in-plane}}
        \;\underbrace{-\,2\left(\Omega_zv_x - \Omega_xv_z\right)\hat{y}_g}_{\text{out-of-plane}}
\end{equation}
This is worth reading carefully, because it inverts the intuitive move. The in-plane Coriolis
acceleration is produced \emph{entirely} by $\Omega_y$ --- the component of $\vec{\Omega}$
\emph{normal} to the flight plane --- while $\Omega_x$ and $\Omega_z$, the components lying
\emph{inside} the plane, produce a purely out-of-plane acceleration. "Project $\vec{\Omega}$ onto
the flight plane and keep that part" would therefore keep exactly the half that does nothing
in-plane and throw away the half that does everything.

\textbf{Simplification (2D closure).} Only the in-plane bracket of
\autoref{eq:coriolis_split} is retained; the $\hat{y}_g$ term is dropped. This is \emph{not} a
smallness argument --- the discarded lateral acceleration is of the same order as the retained
one, and for some geometries larger. It is dropped because the model has no lateral degree of
freedom to receive it: the out-of-plane Coriolis deflection (the familiar right-hand deflection
in the northern hemisphere) is physically real, and a 2D planar model simply cannot represent it.
This is the same closure already applied to the crosswind component $w_y$ in \S Wind, and it is
lifted by the same future work (Phase 3, full 6-DOF), not by anything in this section.

Two consequences follow immediately and should be stated so results are not over-read:
\begin{itemize}
    \item $\Omega_y = -\Omega\cos\phi_0\sin\chi$ vanishes for a launch due North or due South
    ($\chi = 0^\circ, 180^\circ$) and at the poles. For those cases the 2D model predicts
    \emph{zero} Coriolis effect, while the true Coriolis effect is maximal and entirely lateral.
    \item $|\Omega_y|$ is largest for an East/West launch near the equator, where the effect the
    2D model \emph{can} see is also the largest it ever is.
\end{itemize}

\paragraph{In-plane Coriolis term.} Substituting $\Omega_y$ into the retained bracket:
\begin{equation} \label{eq:coriolis_2d}
    \vec{a}_{Coriolis}^{\,2D} = 2\Omega\cos\phi_0\sin\chi\left(v_z\,\hat{x}_g - v_x\,\hat{z}_g\right)
\end{equation}
i.e.\ componentwise, with the state variables of \S Equations of motion ($v_x = \dot{x}$
downrange, $v_z = \dot{z}$ \emph{downward}, so an ascending rocket has $v_z < 0$):
\begin{align}
    a_{Coriolis,x} &= +2\Omega\cos\phi_0\sin\chi\;\dot{z} \\
    a_{Coriolis,z} &= -2\Omega\cos\phi_0\sin\chi\;\dot{x}
\end{align}
The two classical checks both come out right. A rocket rising vertically at the equator with an
Eastward $+\hat{x}_g$ ($\phi_0=0$, $\chi=90^\circ$, $\dot{z}<0$) gets $a_x < 0$: it drifts
\emph{West}, the standard counterpart of the eastward deflection of a falling body. A body moving
East ($\dot{x}>0$) gets $a_z < 0$, i.e.\ \emph{upward} --- the E\"otv\"os effect, eastward motion
reducing apparent weight.

\paragraph{Residual centrifugal term.} Write the rocket's geocentric position as
$\vec{r} = \vec{R}_0 + \vec{s}$, with $\vec{R}_0$ from the Earth's centre to the launch site and
$\vec{s} = x\hat{x}_g + z\hat{z}_g$ the simulation's own position state. Since
$-\vec{\Omega}\times(\vec{\Omega}\times\vec{a}) = \Omega^2\vec{a} - (\vec{\Omega}\cdot\vec{a})
\vec{\Omega}$ is linear in $\vec{a}$, the centrifugal acceleration splits cleanly:
\begin{equation}
    \vec{a}_{centrifugal} = \underbrace{\Omega^2\vec{R}_0 - (\vec{\Omega}\cdot\vec{R}_0)\vec{\Omega}}_{\text{constant, already in }\gamma_0(\phi)}
    \;+\; \underbrace{\Omega^2\vec{s} - (\vec{\Omega}\cdot\vec{s})\vec{\Omega}}_{\text{residual}}
\end{equation}
The first bracket is the static launch-site term the WGS84 Somigliana formula already carries.
The residual, with $\vec{\Omega}\cdot\vec{s} = \Omega(\cos\phi_0\cos\chi\,x - \sin\phi_0\,z)$
from \autoref{eq:omega_components}, has in-plane components
\begin{align}
    a_{cf,x} &= \Omega^2\Bigl[x - \cos\phi_0\cos\chi\left(\cos\phi_0\cos\chi\,x - \sin\phi_0\,z\right)\Bigr] \\
    a_{cf,z} &= \Omega^2\Bigl[z + \sin\phi_0\left(\cos\phi_0\cos\chi\,x - \sin\phi_0\,z\right)\Bigr]
\end{align}
This is bounded by $\Omega^2|\vec{s}| = 5.3\times10^{-9}\,|\vec{s}|$, i.e.\
$5\times10^{-5}\,\mathrm{m/s^2}$ even at $|\vec{s}| = 10\,\mathrm{km}$ --- roughly three orders of
magnitude below the Coriolis term, and its \emph{vertical} part is in any case already folded into
the WGS84 free-air correction through the constant $m$ in \S Gravity (double-counting it would be
a small error in the wrong direction). Recommendation: implement the Coriolis term only, and leave
the residual centrifugal term documented here but unimplemented.

\paragraph{Moment.} Coriolis and centrifugal accelerations are body forces distributed in exact
proportion to mass, so their resultant acts at the centre of gravity $B$ and contributes
\emph{zero} moment about $B$. The one surviving coupling comes from the body's own pitch rate
($\vec{v}_i = \vec{v}_B + \vec{\omega}\times\vec{r}_i$ inside the rigid body), a gyroscopic torque
of order $2I\Omega\omega$, i.e.\ an angular acceleration of order
$2\Omega\omega \approx 1.5\times10^{-4}\,\mathrm{rad/s^2}$ at $\omega = 1\,\mathrm{rad/s}$. It is
negligible against the aerodynamic pitching moment and is dropped; the rotational equation of
\S Equations of motion is unchanged.

\paragraph{Result to add to \S Equations of motion.} The translational equation gains one term:
\begin{equation}
    m\vec{\ddot{x}} = \vec{W} + \vec{T} + \vec{A}_f + m\,\vec{a}_{Coriolis}^{\,2D}
\end{equation}
with $\vec{a}_{Coriolis}^{\,2D}$ from \autoref{eq:coriolis_2d}, evaluated directly in
$\mathcal{F}_0$ (no body-axis rotation needed --- it is already a ground-frame expression, like
$\vec{W}$). $\phi_0$ and $\chi$ are the launch-site latitude and heading already used by
\S Gravity and \S Wind respectively; no new parameter is introduced.

\paragraph{Magnitudes, and an honest consistency caveat.} The retained term scales as
$2\Omega\cos\phi_0\sin\chi\,V$: at $V = 300\,\mathrm{m/s}$ it is at most
$4.4\times10^{-2}\,\mathrm{m/s^2}$ (equatorial East launch) and about
$2.8\times10^{-2}\,\mathrm{m/s^2}$ at $\phi_0 = 50^\circ$ --- some $0.3\%$ of $g$, accumulating to
order $10\,\mathrm{m}$ of displacement over a $30\,\mathrm{s}$ powered ascent. That is worth
having for impact-point prediction and is consistent with the centimetre-to-metre estimate quoted
in \S Current assumptions.

It is, however, the \emph{same} order as an error this model still carries: freezing $\hat{z}_g$ at
the launch-site vertical (\S Gravity > Direction) mis-points gravity by $x/R_e$, an effective
$g\,x/R_e \approx 1.5\times10^{-2}\,\mathrm{m/s^2}$ at $x = 10\,\mathrm{km}$ and growing linearly
with downrange. Adding Coriolis while keeping the frozen-vertical assumption therefore does not
by itself make the trajectory accurate to the Coriolis level; the two corrections are independent
and comparable, and the flat-Earth one dominates once the flight goes far enough downrange. Both
are listed in \S Current assumptions for that reason.}

\subsection{Thrust}
Magnitude of thrust is the result of conservation of momentum (ejection of fuel) and pressure differences, given as
\begin{equation}
    T = \dot{m}v_e + A_e(p_e - p_{atm}),
\end{equation}
which can be computed analytically (internal combustion, nozzle).

In real life, however the data from static test is used. It is given as a tabulated curve,
linearly interpolated between measured $(t_i, T_i)$ points:
\ai{\begin{equation}
    T_{ref}(t) = \mathrm{lerp}\!\left(t; \{t_i\}, \{T_i\}\right)
\end{equation}
held at $T_1$ for $t < t_1$ and zero for $t > t_N$; burnout is the falling zero-crossing of
$T_{ref}(t)$.}

As the test is static, there is no variation in atmospheric pressure. To account for the change in
altitude, \ai{the net thrust once the motor is flowing is}
\begin{equation}
    T(t, h) = T_{ref}(t) + A_e\left(p_{ref} - p_{amb}(h)\right), \qquad \text{only while } T_{ref}(t) > 0
\end{equation}
with $p_{amb}(h)$ from \S Atmosphere, $A_e$ the nozzle exit area, and \ai{$p_{ref}$ the ambient
pressure the curve was measured at (the test-stand pressure --- not necessarily the launch site's
own pad pressure)}; $A_e = 0$ recovers $T_{ref}(t)$ exactly.

\ai{Propellant mass flow and the resulting CoG/moment-of-inertia travel are modelled in \S Mass
and moment of inertia; \S Current assumptions still holds today (constant $m$) until that model is
wired into \S Equations of motion.}

\subsection{Mass and moment of inertia}
\ai{As propellant burns, mass leaves the rocket through the nozzle and both the center of gravity
and the moment of inertia change over the burn. A minimal model, treating the rocket as two lumped
masses --- dry structure (mass $m_d$, fixed centroid at $l_d$ from the nose tip $A$) and remaining
propellant (mass $m_p(t) = m(t) - m_d$, centroid at $l_p$):
\begin{gather}
    \dot{m}(t) = -\frac{T_{ref}(t)}{I_{sp}\,g_0} \qquad \text{(or a measured mass-flow table, if available)} \\
    l_B(t) = \frac{m_d\,l_d + m_p(t)\,l_p}{m_d + m_p(t)} \qquad \text{(CoG distance from the nose,
    replaces the constant $l_B$ of \S Geometry)} \\
    I(t) = I_d + m_d\bigl(l_B(t)-l_d\bigr)^2 + I_p(t) + m_p(t)\bigl(l_B(t)-l_p\bigr)^2 \qquad
    \text{(parallel-axis theorem)}
\end{gather}
with $I_d$, $I_p(t)$ the dry structure's and remaining propellant's own moments of inertia about
their respective centroids, and $I_{sp}$ the motor's specific impulse.
\todo{$I_p(t)$ and $l_p$ depend on the propellant grain geometry (cylindrical, tubular, end-burner,
etc.) and how it recedes as it burns --- needs the motor's grain geometry to go beyond the
"propellant centroid fixed at $l_p$" approximation above. Until resolved, $l_p$ constant and
$I_p(t) \propto m_p(t)$ (uniform density, fixed shape) is the simplest placeholder.}}

\subsection{Aerodynamics}
\paragraph{Force magnitudes.}
\ai{$C_A(\alpha,M)$, $C_N(\alpha,M)$, $C_m(\alpha,M)$ and the CoP-to-CoG distance $l_1(\alpha,M)$
come from one rectangular $(M,\alpha)$ table, or fall back to constants:
\begin{equation}
    C_A = C_{A0}, \qquad C_N = C_{N\alpha}\alpha, \qquad C_m = C_{m\alpha}\alpha, \qquad l_1 = l_1^{(0)}
\end{equation}
With a table: bracket and clamp $M$ at the grid edges (no extrapolation beyond the table's Mach
range), interpolate linearly in $\alpha$ within each bracketing row, then blend the two rows
linearly across the $M$ bracket. $l_1(\alpha,M)$ feeds directly into \S Geometry's
$\vec{aA}=-l_1\hat{x}_r$, coupling CoP travel back into the pitching moment.

\ai{Implementation note: the fallback $C_N$, $C_m$ are evaluated as $C_{N\alpha}\sin\alpha$,
$C_{m\alpha}\sin\alpha$ in code (identical to the formula above for small $\alpha$, within 0.4\%
at $8.6^\circ$) rather than the raw linear form. This is a numerical-robustness fix, not a
different physical model: $\alpha=\operatorname{atan2}(w,u)$ is discontinuous at its own branch
cut ($u<0$, $w$ crossing zero), which an unperturbed flight coasting tail-first through $V=0$ at
apogee lands on exactly, and the unbounded linear form turns that into a solver-breaking jump in
the moment every evaluation. $\sin\alpha$ removes the discontinuity (and, at large $\alpha$, is
arguably more correct besides: it is exactly zero at $\alpha=\pi$, matching the zero net side
force a body of revolution has flying exactly tail-first by symmetry).}

The resulting force is resolved along the body axes directly --- \emph{not} rotated into the
relative-wind frame $\hat{u},\hat{n}$ --- since $C_A, C_N$ are themselves body-axis (axial/normal)
coefficients, the same convention the CFD/wind-tunnel data is extracted in, and their
$\alpha$-dependence is already carried by the table:
\begin{equation}
    \vec{A}_f = -A_x\hat{x}_r - A_z\hat{z}_r, \qquad A_x = qSC_A(\alpha,M), \quad A_z = qSC_N(\alpha,M)
\end{equation}
matching \S Forces' original Phase 1 form exactly, for every $\alpha$ --- not just $\alpha=0$.}

\paragraph{Coefficients.}
\ai{What affects $C_A$, $C_N$, $C_m$: angle of attack $\alpha$ (nonlinear at high $\alpha$ once
flow separates), Mach number $M$ (compressibility, transonic/supersonic drag rise), Reynolds
number $Re$ (boundary-layer state, viscous/skin-friction contribution to $C_A$), and geometry
(fin planform and cant, body-fin interference, surface finish). \todo{expand once the CFD/data
extraction pipeline is defined.}}

\subsection{Atmosphere}
Goal: given altitude $h$ (and, for the interpolated model, horizontal position), provide the
local air density $\rho$, together with everything needed to form the Reynolds and Mach numbers
of the relative flow (\S Relative wind and angle of attack), and the local speed of sound $a$ and
dynamic viscosity $\mu$.

\subsubsection{models}
\ai{Three tiers, increasing in fidelity, all exposing the same $(\rho, T, p, a, \mu)$ interface:
\begin{enumerate}
    \item \textbf{ISA} (International Standard Atmosphere): the standard piecewise-linear
    temperature-lapse model (troposphere + lower stratosphere), pressure from the hydrostatic
    equation, density from the ideal gas law, speed of sound $a = \sqrt{\gamma R_{specific} T}$,
    viscosity from Sutherland's law.
    \item \textbf{Real atmosphere} (single station): measured/forecast sounding data ($T(h)$,
    $p(h)$ at one geographic location and time), interpolated in altitude only.
    \item \textbf{Real atmosphere} (multi-station, spatially interpolated): the same kind of
    measured data from several geographic locations, interpolated in altitude, horizontal
    position, and optionally time, to estimate local conditions along the actual ground track.
\end{enumerate}

\paragraph{ISA.} Geopotential altitude from geometric altitude above mean sea level $h_{msl}$:
\begin{equation}
    H = \frac{R_{isa}h_{msl}}{R_{isa}+h_{msl}}, \qquad R_{isa} = 6356766\,\mathrm{m}
\end{equation}
Seven layers, base geopotential altitude $H_b$, lapse rate $L$, base temperature/pressure $T_b, p_b$:
\begin{center}
\begin{tabular}{l c c c c}
Layer & $H_b$ [m] & $L$ [K/m] & $T_b$ [K] & $p_b$ [Pa] \\
troposphere  & 0     & $-0.0065$ & 288.15 & 101325   \\
tropopause   & 11000 & $0$       & 216.65 & 22632.06 \\
stratosphere & 20000 & $+0.0010$ & 216.65 & 5474.889 \\
stratosphere & 32000 & $+0.0028$ & 228.65 & 868.0187 \\
stratopause  & 47000 & $0$       & 270.65 & 110.9063 \\
mesosphere   & 51000 & $-0.0028$ & 270.65 & 66.93887 \\
mesosphere   & 71000 & $-0.0020$ & 214.65 & 3.956420 \\
\end{tabular}
\end{center}
Within the layer containing $H$:
\begin{gather}
    T = T_b + L(H-H_b) \\
    p = p_b(T/T_b)^{-g_{isa}/(LR_d)} \ \ (L\neq0), \qquad
    p = p_b\exp\!\left(\dfrac{-g_{isa}(H-H_b)}{R_dT_b}\right) \ \ (L=0) \\
    \rho = \frac{p}{R_dT}, \qquad a = \sqrt{\gamma R_dT}, \qquad
    \mu = \mu_0\left(\frac{T}{T_0}\right)^{3/2}\frac{T_0+S_\mu}{T+S_\mu}
\end{gather}
with $R_d = 287.0528\,\mathrm{J/(kg\,K)}$, $\gamma = 1.4$, $g_{isa} = 9.80665\,\mathrm{m/s^2}$ (the
ISA's own fixed value, independent of \S Gravity), $\mu_0 = 1.716\times10^{-5}\,\mathrm{Pa\,s}$,
$T_0 = 273.15\,\mathrm{K}$, $S_\mu = 110.4\,\mathrm{K}$ (Sutherland's law).

\paragraph{Sounding.} One station's measured $(h_{msl}, p, T, u, v)$ profile:
\begin{equation}
    T(h) = \mathrm{lerp}(h;\{h_i\},\{T_i\}), \quad
    p(h) = \exp\mathrm{lerp}(h;\{h_i\},\{\ln p_i\}), \quad
    u(h),v(h) = \mathrm{lerp}(h;\{h_i\},\{u_i,v_i\})
\end{equation}
$\rho, a, \mu$ are derived from the interpolated $T, p$ exactly as ISA above -- never interpolated
directly, to stay thermodynamically consistent. Below the lowest level the profile continues with
the sounding's own local lapse rate; above the highest, an ISA-shaped profile re-anchored at
$(h_{top}, T_{top}, p_{top})$ keeps $T, p$ continuous at the seam. Wind is held at the edge value
outside the data range.

\paragraph{Field.} $N$ stations projected onto a tangent plane centred at the launch site
$(\phi_0,\lambda_0)$, together with the rocket's own ground position $(x\sin\chi, x\cos\chi)$
($x$ = downrange coordinate, $\chi$ = launch heading, degrees clockwise from North, the direction
of $+\hat{x}$):
\begin{equation}
    E = R_e\cos\phi_0\,(\lambda-\lambda_0), \qquad N = R_e\,(\phi-\phi_0)
\end{equation}
weighted by inverse squared distance and combined into one profile:
\begin{equation}
    w_i = \frac{1}{d_i^2+\varepsilon^2}, \qquad \hat{w}_i = \frac{w_i}{\sum_j w_j}, \qquad
    (T, \ln p, u, v) = \sum_i \hat{w}_i\,(T_i, \ln p_i, u_i, v_i)
\end{equation}
then interpolated in altitude exactly as the sounding tier. A single station reduces every
$\hat{w}_i$ to 1 and reproduces the sounding tier exactly.}

data given

\subsubsection{density}
Humid air density calculation
\ai{\todo{not modelled: radiosonde soundings carry dewpoint, and moist air is $\sim$1\% less
dense than dry air at the same $p,T$ (a $T_{virtual}$ correction). A few lines once the dewpoint
column is carried through the sounding loader; deferred to the backlog until then.}}

\subsubsection{reynolds, mach}
\ai{\begin{equation}
    Re = \frac{\rho V l_{ref}}{\mu}, \qquad M = \frac{V}{a}
\end{equation}
with $V$ from \S Relative wind and angle of attack, $l_{ref}$ a reference length (e.g.\ body
diameter), $\mu$ the dynamic viscosity, and $a$ the local speed of sound, both from the tier
above.}

\subsection{Wind}
\ai{Goal: given time $t$ and position (altitude $h$, and horizontal position for the interpolated
tier), provide the wind velocity $\vec{v}_w$ in $\mathcal{F}_0$, used in
$\vec{v}_{rel} = \vec{v} - \vec{v}_w$ (\S Relative wind and angle of attack). Same three-tier
structure as \S Atmosphere:
\begin{enumerate}
    \item \textbf{Simple analytic model}: e.g. constant wind plus an altitude-dependent shear
    profile, optionally with seeded gusts.
    \item \textbf{Real wind data} (single station): measured/forecast wind profile $\vec{v}_w(h)$
    at one location and time.
    \item \textbf{Real wind data} (multi-station, spatially interpolated): wind profiles from
    multiple locations, interpolated in altitude, horizontal position, and optionally time.
\end{enumerate}

\paragraph{Flight-plane projection.} The launch heading $\chi$ (\S Atmosphere) fixes the
2D flight plane. A horizontal wind given as ENU components $(u,v)$, or as meteorological
speed/direction $(W,\psi)$ with $u=-W\sin\psi,\ v=-W\cos\psi$ ($\psi$ = the direction the wind
blows \emph{from}), projects onto it as
\begin{gather}
    w_x = u\sin\chi + v\cos\chi = -W\cos(\psi-\chi) \\
    w_y = u\cos\chi - v\sin\chi = -W\sin(\psi-\chi)
\end{gather}
with $w_x$ the downrange (along-plane) component feeding $\vec{v}_w$ and $w_z = 0$ (no vertical
wind modelled). $w_y$, the cross-plane component, is computed and discarded: crosswind
weathercocking needs the out-of-plane freedom only Phase 3 has.

\paragraph{Analytic.} Power-law shear, constant direction, and optional deterministic gusts:
\begin{gather}
    W(h) = W_{ref}\left(\max(h,h_{ref})/h_{ref}\right)^{n}, \qquad \psi = \psi_0 \\
    \Delta W(t) = \sigma\sum_i w_i\sin(\omega_it+\varphi_i)
\end{gather}
$h$ is altitude above the local ground, $n\approx1/7$ the shear exponent (the classic open-terrain
value; smoother terrain is lower, rougher/urban/forested terrain higher, up to $\sim$0.4), valid
within the atmospheric surface/boundary layer only --- roughly the lowest few hundred meters to
1--2\,km depending on stability, above which the sounding tiers below should be used instead.
$\Delta W(t)$ added to $W(h)$ before the projection above (then clamped non-negative, since $W$
is a speed). The $\{\omega_i,\varphi_i,w_i\}$ are drawn once and frozen -- the adaptive solver
samples the right-hand side at arbitrary, repeated and rejected times, so a fresh random draw per
call would make $\Delta W$ a non-function of $t$. $\sigma=0$ disables gusts entirely.

\paragraph{Sounding / field.} The same measured $(u,v)$ profiles and interpolation as
\S Atmosphere, projected through the formula above; never interpolated as $(W,\psi)$,
since direction wraps at $360^\circ$.

\paragraph{Vertical wind.} All tiers above return horizontal wind only ($w_z = 0$);
updrafts/downdrafts are not modelled (\S Current assumptions). Two standard families exist in the
literature for deriving a vertical component: (1) terrain-slope / mass-continuity orographic lift,
$w \approx \vec{v}_h\cdot\nabla h_{terrain}$, which needs a terrain elevation model this project
does not have; (2) thermal/convective updraft parameterizations (as used in the soaring/glider
literature), which need surface heat-flux or insolation data outside the radiosonde soundings used
here. Neither fits the current data pipeline (station soundings + a simple analytic profile)
without adding a new data source, so this stays unimplemented for now.
\todo{revisit if/when a terrain or NWP vertical-velocity data source is added.}}


\newpage
\ai{\section{Backlog and future work}\label{sec:backlog}
The complement of \autoref{sec:assumptions}: that section records what the present model
simplifies, this one what is deliberately left undone. Nothing here is scheduled. Each entry
says what the item is and --- more usefully --- what has to exist before it can be written down
properly, so that a reader can tell a missing derivation from a missing input.

\subsection*{Modelling and analysis extensions}
\begin{itemize}
    \item \textbf{Thrust-axis misalignment and control.} Give the thrust its own axis, distinct
    from the body axis $\hat{x}_r$: a fixed nozzle misalignment, a commanded gimbal angle, or
    both. $\vec{T}$ then acquires a $\hat{z}_r$ component, the term $\vec{BC}\times\vec{T}$ of
    \autoref{eq:eom_rot} stops vanishing identically, and the thrust point's axial position
    $l_C$ --- irrelevant today, \autoref{sec:analysis_moment} --- becomes a live
    quantity for the first time. This is also the entry point for any control law, a gimbal
    being the only pitch authority a rocket has once it is off the rail. \emph{Waiting on:} a
    defined thrust-axis geometry --- how the deflection is parameterised, and whether it is a
    fixed manufacturing error or a commanded input.
    \item \textbf{Wing-body aerodynamics: analytic CoP, airframe, air brake.} Replace the
    tabulated or constant coefficients of \autoref{sec:aero} with a wing-body
    (Barrowman-style) build-up: nose, body and fins each contributing to $C_N$ and to the
    centre-of-pressure location, so that $l_1(\alpha, M)$ is \emph{computed} from the airframe
    rather than read from a table, and an air brake enters as a switchable change to $C_A$ and
    to the CoP. \emph{Waiting on:} an actual airframe --- nose, body and fin dimensions --- which
    this document does not yet fix anywhere.
    \item \textbf{Dimensionless stability margin.} \autoref{eq:static_margin} is dimensional.
    The rocketry convention is to quote static margin in calibers, i.e.\ divided by the body
    diameter (or by $\sqrt{S}$ for a non-circular reference), which makes the number comparable
    between airframes and removes its dependence on the choice of reference length $d_{ref}$.
    The single reference-area/length convention this was waiting on is now settled --- $S$ and
    $d_{ref}$ in both \autoref{sec:analysis} and \autoref{sec:aero} --- so what is left is the
    rewrite itself, plus one small judgement: whether to report $\textrm{SM}$ or
    $-\textrm{SM}$, so that a stable rocket comes out with a \emph{positive} margin as the
    caliber convention expects.
    \item \textbf{Dynamic stability analysis.} \autoref{sec:analysis} is
    \emph{static}: it asks only whether a disturbance in $\alpha$ produces a restoring moment,
    not what happens next. The dynamic question --- the natural frequency of the resulting pitch
    oscillation, and whether it decays or grows over a flight whose $q$, $M$, $m$ and $I$ are
    all changing --- needs a linearised short-period model built on \autoref{eq:moment_nothrust}.
    \emph{Waiting on:} a pitch-damping derivative $c_{m,q}$, which neither the coefficient table
    nor \autoref{sec:aero} currently provides; without it the analysis can only be
    qualitative. \todo{needs $c_{m,q}$ before it can be quantitative.}
\end{itemize}

\subsection*{Deferrals already stated elsewhere}
Listed here only so that one place in the document answers ``what is not modelled''. Each is
stated in full, with its reasoning, at the point where it is deferred; none of it is repeated
here.
\begin{itemize}
    \item Launch rail constraint, and recovery / parachute --- \autoref{sec:assumptions}.
    \item Gravity \emph{direction} varying with downrange distance, i.e.\ lifting the frozen
    $\hat{z}_g$ --- \autoref{sec:gravity} > Direction. The one genuinely new physics among the
    gravity variants.
    \item Out-of-plane motion: crosswind weathercocking (\autoref{sec:wind}) and the lateral
    Coriolis deflection (\autoref{sec:noninertial}). Both are discarded by the
    same 2D closure and both are lifted by the same future work, a full 6-DOF model --- not by
    anything that can be done inside this document.
    \item Propellant grain geometry behind $I_p(t)$ and $l_p$ ---
    \autoref{sec:mass}.
    \item Humid-air density, i.e.\ the virtual-temperature correction ---
    \autoref{sec:density}.
    \item Vertical wind, updrafts and downdrafts --- \autoref{sec:wind} > Vertical wind.
    \item Expansion of what drives $C_A$, $C_N$, $C_m$ beyond the qualitative list ---
    \autoref{sec:aero} > Coefficients, pending the CFD / data-extraction pipeline.
\end{itemize}}

\end{document}